\chapter{Proofs for Vaught's Formula-Algebra Construction}
\label{annex:p03-n-a-003-01}

\begin{annexitem}{annex:p3.vaught.overview-prime}{thm:p03-prime-atomic}{The prime--atomic overview with its full language and cardinal scope}
\begin{proof}
For a complete theory with infinite models in a relational/constant
language, \cref{thm:p03-vaught-prime-atomic} proves the first clause,
and \cref{thm:p03-vaught-isolation-density} proves the second.
The identification of formula atoms with isolated complete types is
\cref{thm:p03-vaught-algebra-types}. For a countable language with
functions, \cref{thm:p03-vaught-functional-bridge} transports the
models, elementary embeddings and type-space isolation in both directions.
If the complete theory instead has a finite model, the finite case of
\cref{thm:p03-vaught-finite-boundary} makes every model isomorphic to it,
prime and atomic, and every finite-tuple type isolated. Each type space
therefore has its isolated points dense. These alternatives exhaust
satisfiable complete theories and establish both overview clauses.
\end{proof}
\end{annexitem}

\begin{annexitem}{annex:p3.vaught.overview-saturated}{thm:p03-countable-saturated}{The smallness and uniqueness overview}
\begin{proof}
\Cref{thm:p03-vaught-saturation-parameters} identifies the source's
saturation with realization over finite parameter sets. Under that
identification, \cref{thm:p03-vaught-smallness-equivalences} gives a
countably infinite saturated model exactly when every empty-parameter
finite-tuple type space is countable. The full existence construction is
\cref{thm:p03-vaught-saturated-existence}, and uniqueness follows from
\cref{thm:p03-vaught-saturated-uniqueness}. The truth-preserving graph
translation in \cref{thm:p03-vaught-functional-bridge} transfers these
claims to every countable function language. The overview assumes an
infinite model, so no finite-cardinality exception is needed.
\end{proof}
\end{annexitem}

\begin{annexitem}{annex:p3.vaught.algebra-types}{thm:p03-vaught-algebra-types}{The complete syntactic and semantic identifications}
\begin{proof}
Validity modulo \(T\) respects every Boolean operation, so the Boolean
identities descend to the quotient. Its order is
\([\alpha]\leq[\beta]\) exactly when \(\vdash_T\alpha\to\beta\).
A formula is nonzero exactly when its existential closure is consistent:
failure of such consistency means every \(T\)-model satisfies its
universal negation. These observations identify the stated formula
definitions with the algebra definitions.

In a proper dual prime ideal, a contradiction cannot occur, because
upward closure would then include every formula. If both \(\varphi\)
and its negation occurred their conjunction would be such a contradiction.
The decision property therefore chooses exactly one. Every finite
conjunction from \(P\) belongs to \(P\) and is consistent. Replacing its
variables by the new constants makes a satisfiable finite part of
\(\overline P\): interpret the constants by a realizing tuple.
Compactness, \cref{thm:fol-compactness}, realizes the whole set in a
named expansion. Conversely, a consistent complete type is closed
under conjunction and implication and decides every formula, so it
is a dual prime ideal. The realized tuple's truth set equals \(P\):
any formula outside \(P\) has its negation inside \(P\).

Every sentence using the new constants is, after renaming its bound
variables away from \(v_0,\ldots,v_{n-1}\), a constant instance of a
formula in \(F_n(T)\). Replacement of the constants by those free
variables and back is inverse up to bound renaming. This can be
checked on atomic formulas and then through each connective and
quantifier; the fresh bound names prevent capture at the quantifier
step. Consequently the decision property makes \(\overline P\)
complete, and the named expansion of a realizing tuple is a model
of that theory. Two tuples have the same complete type exactly when
their named expansions satisfy the same sentences.

A consistent atom decides each formula. Indeed, if its conjunction
with \(\varphi\) is consistent it entails \(\varphi\); otherwise it
entails \(\neg\varphi\). Its consequences therefore form a complete
type. Conversely, if a consistent formula isolates a complete type,
and it did not decide some \(\varphi\), its two consistent refinements
with \(\varphi\) and \(\neg\varphi\) would extend, by compactness,
to distinct complete types. This contradicts isolation. Thus the
formula is an atom. A principal complete \(P\) contains its own
generator and hence this atom; containing an atom forces \(P\) to
be precisely the atom's consequence set.
For the general principality definition, if a consistent \(\alpha\)
entails every member of a complete \(P\), then \(\alpha\in P\);
otherwise \(\neg\alpha\in P\), contradicting its consistency.
The same argument identifies that definition with this generator.

If a consistent \(\alpha\) has only one expanded elementary theory
among its realizations, it decides every formula: any two consistent
opposite refinements would supply two realizations with different
truth values. Conversely the atom's decision property gives the
same theory to every realization. Consistency is essential here;
the universal implication about realizations alone is vacuous for
an inconsistent formula.

The basic sets \([\varphi]\) have finite intersections
\([\varphi\land\psi]\) and complements \([\neg\varphi]\), and cover
the type space with \([\top]\). A nonempty basic set contains an
isolated point exactly when it contains a generating atom, whose
formula is below that basic formula. This proves atomisticity
equivalent to isolated-point density. Completeness of \(T\) gives
exactly the true and false equivalence classes of sentences in
\(F_0(T)\); the true class generates the unique complete zero-type.
\end{proof}
\end{annexitem}

\begin{annexitem}{annex:p3.vaught.henkin-naming}{thm:p03-vaught-henkin-naming}{The witness step and the named elementary submodel}
\begin{proof}
Take a model of the starting theory and \(\Gamma_s\).
The starting theory has language \(L\), with none of the new \(C\)-symbols.
The chosen \(c_{\nu_s}\) occurs in neither \(\Gamma_s\) nor
\(\varphi_s\), so its interpretation can be changed without affecting
either. If \(\exists v_0\varphi_s\) is true, interpret it by one such
witness; otherwise give it any element of the nonempty domain.
The resulting expansion satisfies \(W_s\) and the preceding conditions.
This proves consistency at the finite stage.

For a sentence with finitely many new constant symbols, first rename
all bound variables away from the free-variable slots to be used.
Replacing the distinct new symbols by the distinct free slots then
gives a formula \(\theta\). A model and interpretations of those
constants satisfy the sentence exactly when the assigned tuple satisfies
\(\theta\); this follows by induction on the formula, since no free
occurrence can be captured. Existentially closing \(\theta\) says
precisely that such interpretations exist. The starting theory contains
none of the new symbols, so it imposes no further restriction on them.
This proves the consistency equivalence, including empty additional
constant lists and repeated assigned values.

Now let a consistent witness theory have a model \(\mathcal N\).
Let \(A=\{c_i^{\mathcal N}:i<\omega\}\).
It is nonempty. Each original individual constant \(d\) has its
value in \(A\): the witness for \(\varphi(v_0)=(v_0=d)\) supplies
a \(c_i\) equal to it. Form the induced relational structure
\(\mathcal M\) on \(A\), with the original constants and all \(C\)
interpreted by their same values. This is an \(L(C)\)-substructure.
For any tuple from \(A\), choose constant names for its entries.
Every existential formula with these parameters is an instance
of one of the enumerated one-free-variable formulas.
If it holds in \(\mathcal N\), its witness axiom supplies a witness in \(A\).

We verify elementarity by induction on formulas. Atomic relations
and equality agree by construction, including nullary relations;
Boolean truth values agree by induction.
For an existential formula, a witness in \(A\) is also a witness in
\(\mathcal N\). In the other direction the witness axiom just described
gives a witness in \(A\), and induction transfers its matrix.
Universals follow by negation, or by transferring a counterexample
existentially. Thus \(\mathcal M\preccurlyeq\mathcal N\) and satisfies
the witness theory. Its reduct has at most countable domain.

For a consistent countable starting theory \(S\), enumerate all
one-free-variable formulas in \(L(C)\) and add just the fresh witness
axioms. Every accumulated finite list is consistent by the first
paragraph. Every finite portion of the resulting theory is therefore
satisfiable, so compactness gives its model; the preceding construction
gives the named countable model. If \(S\) forces infinitude, that
model is countably infinite.

Finally fix any \(c_j\). There are infinitely many syntactically
different formulas
\[
 (v_0=c_j)\land\forall v_{t+1}(v_{t+1}=v_{t+1}),\qquad t<\omega.
\]
Their existential closures are true in every expansion.
Their fresh witnesses therefore force \(c_{\nu_s}=c_j\).
At each such stage this witness symbol occurs in its retained axiom,
so no later such fresh choice can reuse it. Infinitely many distinct
symbols thus name the same value. The argument also works when further
finite conditions are added between the witness steps, provided their
consistency and every witness axiom are retained. It supplies the alias
property used by the simultaneous construction.
\end{proof}
\end{annexitem}

\begin{annexitem}{annex:p3.vaught.simultaneous}{thm:p03-vaught-simultaneous-construction}{Every original Henkin branch and every tuple}
\begin{proof}
Treat \(T_1\) as \(T\) in the language expanded by \(C\), with no
restrictions on the new constants. The original standing assumptions
make its base satisfiable.
The family \(\mathcal P_0(T)\) has no nonprincipal member, so every
omission requirement of a nonprincipal type has positive finite arity.
There are countably many requirements consisting of either an empty
type flag or a listed type, together with a syntactically distinct
tuple of \(C\)-symbols of matching positive length.
List all of them as \(\Pi_s\), and independently list all
one-free-variable formulas as \(\varphi_s\).
Both lists are countably infinite and can be paired by their indices.

Start with \(\Gamma_0=\top\). At stage \(s\), assume \(T_1+\Gamma_s\)
is satisfiable. Add the fresh witness \(W_s\), obtaining the consistent
sentence \(\gamma_s\). Let \(\bar d\) be the current task tuple.
List the other occurring \(C\)-symbols as \(\bar e\).
They are disjoint from \(\bar d\); unused symbols in \(\bar d\)
are harmless. Freshen bound variables, and replace \(\bar d,\bar e\)
by the free slots to obtain \(\theta_s(\bar v,\bar w)\).
Then \(\gamma_s\) is its exact constant instance up to bound renaming.
The projection \(\beta_s=\exists\bar w\,\theta_s\) is consistent
with \(T\), by the projection lemma.

If an atom \(\alpha\) lies below \(\beta_s\), a model and tuple
satisfying \(\alpha\) also satisfy \(\beta_s\). Choose its \(\bar w\)
witness, interpret \(\bar d,\bar e\) by these coordinates and any
other \(C\)-symbols arbitrarily. It satisfies
\(\alpha(\bar d)\land\gamma_s\). Thus the new accumulated sentence
is consistent.
If there is no such atom, \(\beta_s\) is an atomless element by
definition. An empty flag needs no additional condition, so the
consistent \(\gamma_s\) is retained.

For a nonempty nonprincipal task \(P\), a \(\delta\notin P\) with
\(\beta_s\land\delta\) consistent must exist.
Otherwise every \(\delta\notin P\) would be inconsistent with
\(\beta_s\). For each \(\varphi\in P\), \(\neg\varphi\notin P\),
and hence \(\vdash_T\beta_s\to\varphi\).
The consistent \(\beta_s\) would be a principal generator for the
complete \(P\), contradicting nonprincipality.
Choose one such \(\delta\). A tuple satisfying
\(\beta_s\land\delta\) has \(\bar w\) witnesses for \(\theta_s\);
the same constant-interpretation argument makes
\(\delta(\bar d)\land\gamma_s\) satisfiable.
Since \(\gamma_s=\Gamma_s\land W_s\), this is precisely consistency
of the previously accumulated sentences together with the added
\(\delta(\bar d)\land W_s\).
All three source branches therefore retain consistency.

Let \(T_2\) be \(T_1\) with all the chosen stage sentences.
Each finite subset is contained in some consistent stage, so
compactness gives a model. Every enumerated formula has its witness
axiom retained, and the naming lemma gives a countable elementary
named submodel satisfying \(T_2\). Its \(L\)-reduct \(\mathcal M\)
models \(T\), so it is infinite, hence countably infinite.

Take any positive-length tuple \(\bar a\) from \(\mathcal M\).
The infinitely many aliases allow choosing syntactically distinct
constant names \(\bar d\) for its coordinates, even when some
coordinates are equal. At its empty-flag requirement, the atom
branch puts an atom on that tuple; otherwise \(\gamma_s\) entails
\(\beta_s(\bar d)\), since its \(\bar e\)-values witness the projection,
and this \(\beta_s\) is atomless. Later sentences cannot remove those
truths. The empty tuple satisfies the sole zero-type atom.
This proves the first property for every finite arity.

For any listed \(P_j\) and tuple of matching arity, use the same
distinct names at its \(P_j\)-requirement. In the atom branch,
the tuple satisfies an atom \(\alpha\). A nonprincipal complete
\(P_j\) contains no atom, so it contains \(\neg\alpha\) and this tuple
cannot realize it. In the atomless branch, the tuple satisfies
\(\delta\notin P_j\), whereas \(\neg\delta\in P_j\); again it cannot
realize the type. This covers every tuple of every listed arity
and proves simultaneous omission.
Deleting the omission tasks gives the first-only construction;
the full construction already gives the omission-only conclusion.
\end{proof}
\end{annexitem}

\begin{annexitem}{annex:p3.vaught.automatic-omission}{thm:p03-vaught-automatic-omission}{The exact isolated-point closure class}
\begin{proof}
Let \(P\) be nonprincipal and suppose every basic neighborhood
\([\beta]\), \(\beta\in P\), contains an isolated type.
Then every such \(\beta\) has an atom below it and cannot be atomless.
If a tuple in the constructed model realized \(P\), the simultaneous
property would give it either an atom or an atomless formula.
An atom cannot belong to a nonprincipal complete type; an atomless
formula cannot belong to this \(P\). Both possibilities contradict
realization. Thus this entire class of \(P\) is omitted.

To verify the sequence description, enumerate the formulas of \(P\)
as \(\varphi_0,\varphi_1,\ldots\), allowing a finite list to repeat.
For each initial conjunction \(\beta_s\), its neighborhood has
infinitely many isolated points. If it had only finitely many,
with atoms \(\alpha_1,\ldots,\alpha_k\), the nonprincipal \(P\)
would contain each \(\neg\alpha_i\).
The neighborhood of
\(\beta_s\land\neg\alpha_1\land\cdots\land\neg\alpha_k\), still a
member of \(P\), would contain an atom different from the entire
list, a contradiction. Choose successively a new isolated type
in \([\beta_s]\), different from the preceding finitely many choices.
For every \(\varphi\in P\), all sufficiently late choices lie in
\([\varphi]\); for every \(\psi\notin P\), they lie in
\([\neg\psi]\). These clopen neighborhoods form a basis, so the
sequence converges to \(P\).
Conversely, convergence of isolated points to \(P\) puts an
isolated point in every neighborhood eventually. The two descriptions
are therefore equivalent in this countable formula space.
\end{proof}
\end{annexitem}

\begin{annexitem}{annex:p3.vaught.prefix-embedding}{thm:p03-vaught-prefix-embedding}{Identity, constants, every tuple and every formula}
\begin{proof}
If \(a_i=a_j\), choose a prefix containing both and apply its type
agreement to \(v_i=v_j\); then \(b_i=b_j\). Conversely its negation
shows \(a_i\ne a_j\) implies \(b_i\ne b_j\).
The proposed mapping is therefore well defined and injective.
For any original constant \(d\), some \(a_i=d^{\mathcal A}\);
the formula \(v_i=d\) gives \(b_i=d^{\mathcal B}\).
For a relation and any tuple, choose a prefix containing indices
for all its entries. The corresponding atomic formula, with repeated
slots where needed, transfers its truth both ways. The image is thus
an induced relational substructure containing all constant values.

For an arbitrary formula \(\psi\) and finite tuple from \(\mathcal A\),
choose those finitely many indices and a prefix containing them.
Replace the free-variable slots in \(\psi\) by the indicated prefix
slots, first freshening bound variables. The resulting formula belongs
to that prefix's \(F_n(T)\); its truth is exactly the original
tuple truth. Complete-type agreement transfers it to the image
tuple. This verifies every formula, not just atomic ones, and proves
elementarity of the embedding. It also proves that repetitions
anywhere in the enumerating lists do not obstruct the construction.
\end{proof}
\end{annexitem}

\begin{annexitem}{annex:p3.vaught.atomic-parameters}{thm:p03-vaught-atomic-parameters}{An isolator for the concatenated tuple}
\begin{proof}
For any finite \(\bar b\), atomicity supplies an atom
\(\alpha(\bar u,\bar v)\) at the concatenated tuple
\((\bar a,\bar b)\). Substitute the new names for \(\bar a\),
obtaining \(\alpha(\bar a,\bar v)\) in the named language.
It is consistent with \(\operatorname{Th}(\mathcal A,\bar a)\),
since \(\bar b\) realizes it.
Any named-language formula \(\psi(\bar v,\bar a)\) is obtained
by substituting \(\bar a\) into an \(L\)-formula \(\psi(\bar v,\bar u)\),
after bound-variable freshening. The original atom decides this
formula over \(T\). Its substituted implication, positive or negative,
therefore holds in every model of the complete named theory.
Thus the substituted atom isolates the full type of \(\bar b\)
over the finite named tuple. Every finite \(\bar b\) has such an
isolator, proving atomicity of the expansion.
\end{proof}
\end{annexitem}

\begin{annexitem}{annex:p3.vaught.atomic-embedding}{thm:p03-vaught-atomic-embedding}{Every step of the source's forth construction}
\begin{proof}
Start the enumeration \(a_0,a_1,\ldots\) of the source with its
prescribed finite tuple, allowing repetitions there.
Keep the target's prescribed tuple as \(b_0,\ldots,b_{m-1}\).
Suppose the prefixes of length \(n\geq m\) have the same complete
type. Take an atom \(\alpha(\bar v)\) at the source prefix and an
atom \(\alpha'(\bar v,v_n)\) at the source prefix extended by \(a_n\).
For \(n=0\), use the zero-type atom \(\top\).
The conjunction \(\alpha\land\exists v_n\alpha'\) is consistent,
because the actual source prefix and element realize it.
Atomicity of the formula \(\alpha\), rather than of the target,
therefore yields
\[
                  \vdash_T\alpha\to\exists v_n\alpha'.
\]
The target prefix satisfies \(\alpha\) by its type agreement.
Choose a witness \(b_n\) for \(\alpha'\) in \(\mathcal B\).
The extended prefixes both satisfy the same atom, so the
algebra/type lemma gives their full type agreement.
This extends the invariant by one coordinate.
Induction defines all \(b_n\); the prefix-embedding criterion makes
their correspondence an elementary embedding of the whole source.
The prescribed coordinates were fixed at the initial stage.

This also gives the local one-new-element extension form.
An element not in the old finite source set satisfies inequalities
to those entries in its extended type; type agreement forces its
new target value to be outside the old range. Repeated source
elements instead have their existing images forced by equality.
These details supply the extensions used in back-and-forth.
\end{proof}
\end{annexitem}

\begin{annexitem}{annex:p3.vaught.atomic-uniqueness}{thm:p03-vaught-atomic-uniqueness}{Forth and back with all enumeration entries reached}
\begin{proof}
The prescribed equal-type tuples give a well-defined finite elementary
bijection between their sets of entries, by the equality formulas.
Choose enumerations of both countably infinite domains.
At each even stage take the least-index source element outside the
current domain and apply the local extension step from the preceding
proof, using source atomicity and the current elementary tuple map.
Its new image is outside the range.
At each odd stage take the least-index target element outside the
range. Apply the same extension argument to the inverse tuple map,
now using the target's atomicity, and invert the enlarged map.
Every finite stage remains an elementary bijection between finite
sets and extends the preceding one.

Each source element is eventually included: after its finitely many
earlier enumeration entries have been included, the next even stage
that needs it chooses it. The same argument on odd stages covers
every target entry. The union is therefore a bijection.
For any formula and finite tuple in the source, some stage contains
all its entries and already preserves the formula's truth.
The union hence preserves all formulas and is an isomorphism.
It continues to extend the original prescribed tuple map.
\end{proof}
\end{annexitem}

\begin{annexitem}{annex:p3.vaught.atomic-homogeneity}{thm:p03-vaught-atomic-homogeneity}{The prescribed map in the same model}
\begin{proof}
Apply the strengthened uniqueness theorem with the same countably
infinite atomic model on both sides, starting with the two finite
tuples of the same complete type. Its resulting isomorphism from
the model onto itself is an automorphism and carries every prescribed
coordinate to its counterpart. This is exactly the stated homogeneity.
\end{proof}
\end{annexitem}

\begin{annexitem}{annex:p3.vaught.prime-atomic}{thm:p03-vaught-prime-atomic}{The full converse through omission}
\begin{proof}
A countably infinite atomic model embeds elementarily in every
\(T\)-model by the atomic embedding lemma, so it is prime.
Conversely let \(\mathcal A\) be prime. The naming construction supplies
a countably infinite \(T\)-model. An elementary embedding into that
model is injective, so \(\mathcal A\) is at most countable; it is
infinite by the standing hypothesis, hence countably infinite.

If it were not atomic, some finite tuple would satisfy no atom.
Its complete type \(P\), the truth set of that tuple, is nonprincipal
by the algebra/type lemma. Apply the simultaneous construction to
the single listed type \(P\); the resulting countably infinite
\(\mathcal B\models T\) omits it. Any elementary embedding
\(\mathcal A\to\mathcal B\) would preserve every formula in that
tuple's \(P\), giving a realization there, a contradiction.
Thus the prime model is atomic. Uniqueness follows from countability
and the proved atomic back-and-forth theorem.
\end{proof}
\end{annexitem}

\begin{annexitem}{annex:p3.vaught.isolation-density}{thm:p03-vaught-isolation-density}{Every direction and the arbitrary-sized atomic premise}
\begin{proof}
The prime-to-atomic implication was just proved.
Suppose an atomic model \(\mathcal A\), of any cardinality, exists.
For a consistent \(\varphi\in F_n(T)\), the existential closure
is true in some \(T\)-model. Completeness of \(T\) then makes that
sentence true in every \(T\)-model, in particular in \(\mathcal A\).
Choose a realizing tuple. Its atomicity supplies an atom \(\alpha\)
at that tuple. Since \(\alpha\land\varphi\) is consistent, the atom
entails \(\varphi\). Thus every consistent formula has an atom
below it, at every finite arity, including the automatic zero case.

Conversely, if every \(F_n(T)\) is atomistic, use the simultaneous
construction's first property. No atomless formula exists, so every
tuple in that countably infinite model satisfies an atom.
It is atomic and hence prime by the embedding theorem.
Finally, equivalence with density of isolated types is the proved
basic-open algebra/type correspondence. All four directions are
therefore supplied without turning the arbitrary atomic model in
the premise into a claimed prime model of that same cardinality.
\end{proof}
\end{annexitem}

\begin{annexitem}{annex:p3.vaught.nonprincipal-algebra}{thm:p03-vaught-nonprincipal-algebra}{Finite atoms and the nonprincipal ultrafilter construction}
\begin{proof}
In a finite Boolean algebra every nonzero element contains a minimal
nonzero element, an atom. Their finite join is one; otherwise its
nonzero complement would contain another atom.
A dual prime ideal must contain one of these finitely many atoms:
if it contained none, it would contain all their complements and
then their zero intersection. It is therefore principal.

Now let a countable Boolean algebra be infinite.
If it has infinitely many atoms, the complements of all atoms have
the finite intersection property. No finite atom list can cover one
while another atom remains. Thus their finite intersections are
nonzero. Enumerate the algebra, and extend these complements to a
complete filter by successively choosing an element or its complement,
whichever preserves finite nonzero intersections with the existing
filter conditions. At least one choice does: otherwise two finite
obstructions, combined, would make an existing finite intersection
zero. The resulting filter is proper, decides every element, and
contains no atom. It is nonprincipal.

If there are only finitely many atoms, their join cannot be one.
Were it one, every algebra element would be the join of some of
those atoms, giving a finite algebra.
Its nonzero complementary element is atomless.
Start a proper filter with that element and complete it by the same
enumeration procedure. It contains no atom: an atom in the filter
would have nonzero intersection with that atomless element, forcing
the atom to lie below it. The resulting complete filter is again
nonprincipal. The case of no atoms is included, with the starting
element one.
Applying this algebra result to the quotient and pulling the filter
back to formulas gives the source's dual prime ideal.
\end{proof}
\end{annexitem}

\begin{annexitem}{annex:p3.vaught.saturation-parameters}{thm:p03-vaught-saturation-parameters}{Finite named parameters and the separate countable-set obstruction}
\begin{proof}
Fix a finite tuple \(\bar a\) of length \(m\) and its complete type
\(P_{\bar a}\). A complete \(r\)-type over these names is consistent
with \(\operatorname{Th}(\mathcal M,\bar a)\), so compactness realizes
it in a model of that named theory. The concatenated tuple of the
named values and realizing entries has a complete \((m+r)\)-type
over \(T\), extending \(P_{\bar a}\).
This type is determined by the named type: every formula in the
concatenated variables becomes a formula over \(\bar a\) by proper
substitution. Conversely it determines that named type by this same
substitution. Thus the correspondence is injective and has exactly
the compatible extending types as its range.

Suppose the source one-coordinate extension property and weak
saturation hold. For the desired named \(r\)-type, take the
concatenated complete type \(Q\) just described.
Its successive restrictions to \(m,m+1,\ldots,m+r\) coordinates
are complete types. Starting with the actual \(\bar a\), use the
one-coordinate property successively to realize each next restriction.
After finitely many steps its \(r\) new entries realize the entire
named type at the specified parameters. The empty tuple of new
entries needs no step. This proves finite-parameter saturation.

Conversely finite-parameter saturation realizes all types over the
empty set, giving weak saturation. Given \(P\subseteq Q\) and
\(\bar a\in P(\mathcal M)\), form the named one-type
\(\{\varphi(\bar a,y):\varphi(\bar v,v_n)\in Q\}\).
For each finite conjunction \(\delta\) from \(Q\),
\(\exists v_n\delta\) belongs to \(Q\cap F_n(T)=P\):
it is entailed by \(\delta\), hence in the complete filter.
Therefore the actual \(\bar a\) has a witness for that conjunction
in \(\mathcal M\).
Every finite part of the named type together with its named theory
is satisfiable; compactness proves consistency. It decides every
named formula, so is complete. Finite-parameter saturation realizes
it, proving the source extension property, including \(n=0\).

The injective correspondence above bounds finite-parameter spaces
by \(\mathcal P_{m+r}(T)\), giving the countability statement.
For the other parameter scope, name every member of an infinite
countable \(\mathcal M\) and add
\[
                         x\ne a\quad(a\in\dom\mathcal M).
\]
Any finite list is realized in \(\mathcal M\) outside its finite
set of excluded entries. Compactness gives a realizing elementary
extension of the full named diagram; its realized truth set is
a complete one-type over the named model. None of its realizations can be in
\(\mathcal M\), since it excludes every element.
This is a type over a countable parameter set, not over a finite one.
\end{proof}
\end{annexitem}

\begin{annexitem}{annex:p3.vaught.interpretation-preservation}{thm:p03-vaught-interpretation-preservation}{Every translation clause and the union}
\begin{proof}
For a distinguished constant, its unique defining formula selects
one domain element. A fresh existential variable constrained by
that formula in an atomic translation is therefore forced to
that element, including when a constant occurs several times.
Relations agree by their defining \(\rho_R\); equality agrees
because the domain uses actual elements and actual identity.
This proves the atomic truth comparison for domain assignments.
Conjunction, disjunction and negation preserve it by induction.
For an existential, a domain witness on the interpreted side
is exactly a target witness satisfying \(D\) and the translated
matrix. The universal guard quantifies precisely over that same
domain. These cases prove the truth lemma for every formula.
The interpretation requirements then make every denotation a
\(T\)-model, and in particular give it an infinite domain.

If \(\mathcal B\preccurlyeq\mathcal B'\), each old element has
the same truth value for \(D\) in both. The old interpreted domain
is therefore the intersection of the new one with \(\dom\mathcal B\).
Every defining relation agrees on old tuples.
An old unique distinguished element remains the new unique one,
since its defining formula holds in both and uniqueness is a
target-theory condition. Thus these denotations are nested
substructures. For every old domain tuple, elementarity transfers
\(\varphi^I\); the truth lemma on both sides transfers \(\varphi\).
Their inclusion is elementary.

For a chain with elementary union \(\mathcal B_\infty\), each element
of the union lies in some stage. By elementarity it satisfies \(D\)
in the union exactly when it does in that stage.
Consequently \(D(\mathcal B_\infty)\) is the union of the stage
domains. Every finite tuple lies in one stage; the same transfer
identifies all defined relations. The unique distinguished
elements are already present in the first stage and remain fixed.
This proves equality of the interpreted union with the denotation
of the target union, as well as the corresponding elementary chain.
\end{proof}
\end{annexitem}

\begin{annexitem}{annex:p3.vaught.interpreted-extension}{thm:p03-vaught-interpreted-extension}{The full diagram construction and its witness}
\begin{proof}
Add a new individual name \(\dot b\) for every element of
\(\mathcal B\), and a further fresh name \(d\).
The elementary diagram is every sentence with these names true in
\(\mathcal B\), including the negative truths.
Add \(D(d)\) and, for every \(\varphi\in Q\), the translated sentence
\(\varphi^I(\dot{\bar a},d)\).

For finitely many such \(\varphi\), their conjunction belongs to
the complete \(Q\). Its existential closure in the last variable
is in \(Q\cap F_n(T)=P\).
The actual \(\bar a\) thus has a witness \(y\) in
\(\mathcal A=I(\mathcal B)\) satisfying all those matrices.
The truth lemma gives \(D(y)\) and all their translated truths
in \(\mathcal B\). Interpret \(d\) by \(y\) and every \(\dot b\)
by its original element. This satisfies any finite diagram portion
and all the finitely requested translated formulas at once.
The argument also covers an empty list of matrices, since the
interpreted domain is nonempty.
Compactness makes the full expanded theory consistent.

Its language is countable: the target, its element names and \(d\)
are all countable. The named-model lemma gives an at most countable
model \(\mathcal N\) of this theory.
Mapping \(b\) to \(\dot b^{\mathcal N}\) is well defined, injective
and elementary, by all positive and negative diagram sentences.
It respects each original constant by its named equality.
Its image is therefore an elementary copy of \(\mathcal B\),
and \(\mathcal N\) is infinite because that image is infinite.
Replace \(\mathcal N\) by an isomorphic copy whose domain consists
of the original \(\mathcal B\) together with disjointly tagged
copies of the remaining elements, with that elementary copy
identified with \(\mathcal B\).
The result is a countably infinite elementary extension
\(\mathcal B^*\) as an actual inclusion.

The \(d\)-value is in its interpreted domain by the explicit
domain sentence. All translated \(Q\)-formulas hold there.
The truth lemma therefore makes
\((\bar a,d^{\mathcal B^*})\) a realization of \(Q\) in
\(\mathcal A^*\), while interpretation preservation keeps the
old \(\bar a\) and its old truth values.
No completeness assumption on \(T_1\) entered this diagram argument.
\end{proof}
\end{annexitem}

\begin{annexitem}{annex:p3.vaught.saturated-existence}{thm:p03-vaught-saturated-existence}{Both enumerations of the full original Theorem 4.2}
\begin{proof}
A target model is infinite because a nonempty interpretation
family gives it the domain of an infinite \(T\)-model.
The consistent countable \(T_1\) therefore has a countably infinite
model \(\mathcal B_0\) by the naming lemma.

First fix any countably infinite target stage \(\mathcal B\).
For each \(i\in J\), every finite tuple \(\bar a\) from
\(I_i(\mathcal B)\) has a unique complete type \(P\).
For every \(Q\in\mathcal P_{n+1}(T)\) extending that \(P\), place
the task \((i,n,P,Q,\bar a)\) on a list.
There are countably many tasks: \(J\) and the finite arities are
countable, finite tuples in a countable domain are countable,
and every indicated type space is countable.
The empty tuple tasks are included.
Enumerate all these tasks. Starting with \(\mathcal C_0=\mathcal B\),
apply the interpreted-extension lemma at the \(s\)-th task to
produce \(\mathcal C_{s+1}\succeq\mathcal C_s\).
Its old tuple still realizes the same \(P\), by interpretation
preservation; hence every listed application is legitimate even
after earlier extensions. Use actual inclusions as constructed there.

The elementary chain theorem \cref{thm:elementary-chain} gives the
countable union \(\mathcal B^*=\bigcup_s\mathcal C_s\), with
\(\mathcal B\preccurlyeq\mathcal B^*\).
Interpretation commutes with this union.
Every enumerated task has its witness in some stage and retains
it in \(\mathcal B^*\). This operation fills all tasks for tuples
of the original \(\mathcal B\), without claiming it filled tasks
for every new tuple introduced while doing so.

Now repeat this operation in an outer chain:
\(\mathcal B_{m+1}=(\mathcal B_m)^*\).
Again \(\mathcal B_\infty=\bigcup_m\mathcal B_m\) is a countably
infinite elementary \(T_1\)-model.
Take any \(i\), any finite tuple \(\bar a\) in
\(I_i(\mathcal B_\infty)\), and any complete \(Q\) extending its
prefix type \(P\).
Some outer stage \(\mathcal B_m\) contains all entries of \(\bar a\);
elementarity transfers their \(D_i\)-membership and full interpreted
type to that stage. Its task \((i,n,P,Q,\bar a)\) was thus included
when making \(\mathcal B_{m+1}\), so a witness for \(Q\) is in that
stage and remains in the final union.
This proves the one-coordinate extension property for every tuple
of the final interpreted model, not just initial tuples.

Weak saturation also follows. The empty tuple realizes the unique
zero-type. For a specified complete \(r\)-type, take its successive
restrictions to \(0,1,\ldots,r\) slots and apply the proved extension
property finitely many times inside the final model.
The resulting tuple realizes that full type.
Therefore each denotation is saturated in the original sense.
The finite-parameter bridge gives the modern interpretation of
that conclusion. The identity interpretation uses domain \(x=x\),
the original relations and the original constant values; its
denotation is just the target model itself, yielding the stated
identity special case.
\end{proof}
\end{annexitem}

\begin{annexitem}{annex:p3.vaught.saturated-embedding}{thm:p03-vaught-saturated-embedding}{Every source tuple extension into the saturated target}
\begin{proof}
Enumerate the source starting with the given finite tuple.
Assume equal-type prefixes \(\bar a,\bar b\) of length \(n\)
have been matched. Let \(P\) be their common complete type,
and let \(Q\) be the full type of the source prefix extended
by its next element. Then \(Q\cap F_n(T)=P\).
The target's saturation supplies a next \(b_n\) extending its
actual \(\bar b\) to a realization of this \(Q\).
Thus the longer prefixes again have identical complete types.
For an empty initial prefix use the zero-type and the same
extension property. Repeated entries are forced to their existing
images by equality; a new source entry is forced outside the old
range by the inequalities in \(Q\).
Recursion matches the whole countable source, and the full-prefix
criterion proves a well-defined elementary embedding with every
prescribed coordinate fixed.
\end{proof}
\end{annexitem}

\begin{annexitem}{annex:p3.vaught.saturated-uniqueness}{thm:p03-vaught-saturated-uniqueness}{Every forth and back stage without an isolation assumption}
\begin{proof}
Start with the finite elementary bijection supplied by the given
equal-type tuples. At even stages choose the least not-yet-matched
entry in the first model and use the preceding one-coordinate
extension argument in the second, saturated model.
At odd stages choose the least not-yet-matched entry in the
second model and extend the inverse map into the first,
saturated model. The extended complete types force distinct
new entries to distinct new entries. All maps are increasing
finite elementary bijections.

Both countable enumerations are exhausted by these alternating
least-entry choices. Their union is a bijection; every finite
tuple and every formula's truth is already preserved at a stage
containing that tuple. The union is therefore an isomorphism
extending the initial finite correspondence.
The source and target saturation conditions supply the two
directions; no atomic isolator was assumed.
\end{proof}
\end{annexitem}

\begin{annexitem}{annex:p3.vaught.saturated-homogeneity}{thm:p03-vaught-saturated-homogeneity}{The back-and-forth automorphism}
\begin{proof}
Use the saturated uniqueness proof with the same model on both
sides and the two equal-type tuples as its initial correspondence.
Its surjective elementary bijection is the required automorphism.
\end{proof}
\end{annexitem}

\begin{annexitem}{annex:p3.vaught.saturation-characterization}{thm:p03-vaught-saturation-characterization}{Each implication in original Theorem 4.6}
\begin{proof}
A countably infinite saturated target admits the elementary
embedding of every countably infinite \(T\)-model by the
saturated-embedding lemma, so it is \(\aleph_0\)-universal.
Its homogeneity was proved separately. This gives \(1\Rightarrow2\).
For any complete type \(P\), the consistent named theory
\(\overline P\) has a countable model by the naming lemma.
Its reduct is infinite and realizes \(P\).
Universality embeds that reduct in the given model, preserving
the realizing tuple's every formula. Hence a universal model
is weakly saturated, giving \(2\Rightarrow3\).

For \(3\Rightarrow1\), take an actual prefix
\(\bar a\in P(\mathcal A)\) and a complete \(Q\) extending \(P\).
Weak saturation gives some \((\bar a',b')\in Q(\mathcal A)\).
Its prefix realizes \(P\), so the stated homogeneity yields an
automorphism \(f\) sending \(\bar a'\) to \(\bar a\).
Automorphisms preserve every formula in \(Q\); thus
\((\bar a,f(b'))\in Q(\mathcal A)\).
This supplies the exact source extension property at the specified
parameters. Weak saturation is already part of clause 3, so
the model is saturated. The empty-prefix case requires only
its already supplied realization of \(Q\).
\end{proof}
\end{annexitem}

\begin{annexitem}{annex:p3.vaught.countable-types-atomistic}{thm:p03-vaught-countable-types-atomistic}{A distinct complete type on every infinite binary branch}
\begin{proof}
Suppose a nonzero atomless element \(\beta\) existed.
Every nonzero element below it is atomless.
Such an element \(\gamma\) is not an atom, so there is a
Boolean element \(\delta\) for which both
\(\gamma\land\delta\) and \(\gamma\land\neg\delta\) are nonzero.
Otherwise \(\gamma\) would decide every algebra element and
would be an atom.
Starting with \(\chi_{\varnothing}=\beta\), split each node
in this way into two nonzero disjoint children.
Choose the least suitable split from a fixed algebra enumeration
at each node. This produces a tree \(\chi_s\), \(s\in2^{<\omega}\),
with nonzero descending branch chains and disjoint siblings.

For each infinite branch \(\eta\), its descending chain has
the finite intersection property. Complete it to an ultrafilter
by the countable enumeration method used in the preceding
nonprincipal-algebra proof, always selecting the first alternative
that preserves all finite intersections with the branch conditions.
Such an alternative exists by the same two-obstruction argument.
This defines a complete filter \(P_\eta\) for every branch.
Two different branches diverge at disjoint children, and their
filters contain those children, so their filters differ.
Thus \(\eta\mapsto P_\eta\) is an injection from \(2^\omega\)
into \(\mathcal P_n(T)\).
This contradicts countability. Hence no such atomless \(\beta\)
exists, and every nonzero formula contains an atom.
\end{proof}
\end{annexitem}

\begin{annexitem}{annex:p3.vaught.smallness-equivalences}{thm:p03-vaught-smallness-equivalences}{All four directions with their cardinal scopes}
\begin{proof}
If all the empty-parameter type spaces are countable, the full
interpretation existence theorem with the identity interpretation
gives a countably infinite saturated model.
That model is universal by the saturated-embedding theorem.
A universal model is weakly saturated by the realization-and-
embedding argument in the saturation characterization proof.
Finally, a countable weakly saturated model realizes every
complete \(n\)-type with one of its countably many finite \(n\)-tuples.
The map taking a tuple to its complete type is a surjection onto
\(\mathcal P_n(T)\); this space is therefore countable.
For \(n=0\) it has one point already.
These implications close the cycle and prove all equivalences.
The separate uniqueness and homogeneity proofs apply to the
countable saturated model. Every realization over named parameters
in those statements has the finite-parameter scope previously proved.
\end{proof}
\end{annexitem}

\begin{annexitem}{annex:p3.vaught.universal-prime}{thm:p03-vaught-universal-prime}{The atomistic consequence at every arity}
\begin{proof}
A universal model implies smallness by the preceding equivalence.
For every \(n\), countability of \(\mathcal P_n(T)\) implies
atomisticity by the full binary-tree lemma.
The prime-existence theorem therefore gives a prime model.
The argument constructs this existence through the atomic
Henkin model and does not identify that prime model with the
possibly different universal model.

For the source's alternative omission route, smallness makes the
union of all nonprincipal finite-arity type spaces countable.
Omit that entire list by the proved omission construction.
Its countably infinite model can realize no nonprincipal complete
finite-tuple type, since each such type was listed. Every realized
type is therefore principal, making the model atomic and hence prime.
An empty nonprincipal list is covered by the ordinary naming construction.
This route uses omission alone rather than the extra atom/atomless
tuple property or the binary-tree atomisticity lemma.
\end{proof}
\end{annexitem}

\begin{annexitem}{annex:p3.vaught.dlo-example}{thm:p03-vaught-dlo-example}{Every order pattern, the cuts and the two parameter languages}
\begin{proof}
A finite tuple in a linear order determines a partition by equality
and a linear ordering of its distinct-value blocks.
Any two tuples with that same pattern, in any two dense orders
without endpoints, can be extended by one entry on either side:
an equal entry has its forced equal counterpart; a new entry belongs
below all existing values, above all of them, or to an interval
between consecutive distinct tuple values.
No endpoints supplies the first two choices, and density the third.
These witnesses can be chosen different from the old finite entries.

Induction on formulas now shows that same-pattern tuples satisfy
the same formulas. Atomic equality and order agree by hypothesis.
The Boolean cases follow by induction.
For an existential formula, extend the tuple by its source witness,
choose a same-pattern target entry as just described, and transfer
the matrix by the induction hypothesis. The reverse implication
uses the extension in the opposite direction.
Universals follow by negation. With empty tuples, this proves
completeness of DLO; with arbitrary tuples, each pattern isolates
its full type. Every abstract block pattern is realized, for instance
by choosing its finitely many ordered rational values.
Partitioning \(n\) labelled slots into \(k\) blocks and ordering those
blocks yields exactly the displayed sum. The empty pattern has
one realization type and corresponds to \(k=0\).

The rational model is countably infinite and atomic, so is prime
and homogeneous by the already proved package.
Its type spaces are finite, hence small, so a countable saturated
model exists. Every countable DLO model is atomic by the pattern
proof; uniqueness of countable atomic models makes that saturated
model isomorphic to the rational one. Hence the rational one is
also saturated and uniquely countable.

The same formula-induction argument with finitely many fixed
rational parameters proves \(\mathbb Q\preccurlyeq\mathbb R\)
as orders, including in the language naming every rational:
each individual formula uses only finitely many names.
For each irrational real \(r\), take the complete type of \(r\)
over these names in the real model.
It contains \(q<x\) for \(q<r\) and \(x<q\) for \(r<q\).
Different real cuts differ at a rational strictly between them.
There are continuum many such cuts, and at most continuum many
types since the parameter language has countably many formulas.
This proves the exact type count over the countable parameter set
and disproves \(\omega\)-stability of the original small DLO theory.

In the expanded complete \(T_{\mathbb Q}\), every rational tuple
is isolated by the conjunction equating its slots with their names.
Every other formula at that tuple becomes a sentence after
substitution, and completeness decides it; those equalities
therefore isolate the full type. Thus the expanded rational model
is atomic and prime.
Its empty-parameter type space already includes all the displayed
irrational cuts, because those rational names are now part of the
language rather than extra parameters.
Every such type is nonprincipal. If a consistent formula isolated
one, completeness of \(T_{\mathbb Q}\) would make its existential
closure true in the rational model. A rational \(q\) would realize
it there, but the irrational cut type contains \(x\ne c_q\).
This contradicts isolation. The rational model omits all those types.
The smallness equivalence rules out a countable weakly saturated
model of the expanded theory. This also proves that finite naming
and countable naming cannot be interchanged in the smallness bridge.
\end{proof}
\end{annexitem}

\begin{annexitem}{annex:p3.vaught.discrete-example}{thm:p03-vaught-discrete-example}{The defining finite position and the infinite predecessor type}
\begin{proof}
For each \(k\), let \(\gamma_k(x)\) say that there are precisely
\(k\) different elements less than \(x\): use \(k\) existential
variables, their pairwise inequalities, their inequalities below
\(x\), and a universal clause that every predecessor is one of them.
For \(k=0\) this is no-predecessor.
In \((\omega,<)\), exactly the element \(k\) satisfies it.
The sentences asserting existence and uniqueness of that value
belong to \(T_{\rm disc}\).

For any finite tuple of natural numbers \(k_0,\ldots,k_{n-1}\),
the conjunction \(\bigwedge_i\gamma_{k_i}(v_i)\) is consistent.
For any formula \(\varphi(\bar v)\), its implication or negative
implication from that conjunction is true in \(\omega\), according
to the truth at this sole tuple. It therefore belongs to the
complete \(T_{\rm disc}\).
The conjunction is an atom at the tuple. Hence this countable
model is atomic and prime.

Let \(\delta_k(x)\) say that \(x\) has at least \(k\) different
predecessors. Every finite subset of the set of these formulas
is realized by a sufficiently large natural number.
Its union is therefore consistent by compactness.
It is nonprincipal: any consistent alleged generator \(\theta(x)\)
has its existential closure true in \(\omega\), by completeness.
A realizing natural number \(j\) has only \(j\) predecessors,
contradicting the generator's implication of \(\delta_{j+1}\).
A complete extension can be obtained as the truth type of a
realization in a compactness model. If this completion had an
isolator, that isolator would also generate all the \(\delta_k\),
which is impossible. Thus every completion is nonprincipal.
None can be realized in \(\omega\); the prime model therefore is
not weakly saturated. No decision procedure for all discrete-order
types was used in these conclusions.
\end{proof}
\end{annexitem}

\begin{annexitem}{annex:p3.vaught.unary-obstruction}{thm:p03-vaught-unary-obstruction}{The independent predicate patterns and atomless refinements}
\begin{proof}
The axioms say, for each finite disjoint \(A,B\subseteq\omega\)
and each positive integer \(m\), that at least \(m\) different
elements satisfy all \(U_i\) for \(i\in A\) and all \(\neg U_i\)
for \(i\in B\).
A model is \(2^\omega\times\omega\), interpreting \(U_i(s,k)\)
by the \(i\)-th bit of \(s\).
For any fixed finite predicate sublanguage, each Boolean cell
is infinite in every model.
Finite tuples agreeing on equality and those predicate values
extend in either direction: preserve an old equality, or choose
a new element of the required infinite cell outside the finite
old tuple. Formula induction using these extensions, just as in
the DLO proof, transfers every formula of that finite sublanguage.
Every sentence of the full language uses finitely many predicates,
so it has the same truth in all models. This proves completeness.

Take a consistent one-variable formula \(\varphi(x)\), which
mentions only a finite predicate set \(E\), and choose \(j\notin E\).
The relevant finite cell of a realization of \(\varphi\) has
both a \(U_j\)-extension and a \(\neg U_j\)-extension, with infinitely
many elements in each, in the displayed bit model.
They satisfy exactly the same \(E\)-language one-variable formulas,
by the preceding finite-sublanguage induction.
Thus both \(\varphi\land U_j\) and
\(\varphi\land\neg U_j\) are consistent.
In particular no consistent one-variable formula is an atom.
Any formula below a nonzero one-variable formula is likewise
consistent and therefore not an atom, so every such formula is
atomless. The prime/atomic existence equivalence excludes both
sorts of model.

Every bit sequence occurs in the bit model and determines a
complete one-type, since the same finite-sublanguage induction
decides every formula from its finite predicate values.
Distinct sequences differ on a predicate. There are continuum
many types, and no more than continuum many because the full
formula set is countable.
The simultaneous model's unary tuples must therefore use its
atomless alternative here; the theorem did not promise an
atomic model when the formula algebra lacked atoms.
\end{proof}
\end{annexitem}

\begin{annexitem}{annex:p3.vaught.finite-boundary}{thm:p03-vaught-finite-boundary}{The finite signature reducts and the infinite-cardinality counterexample}
\begin{proof}
Let the finite model have size \(k\geq1\).
Its complete theory contains the exact-cardinality sentence,
so every other model also has \(k\) elements.
For a finite collection \(L_0\) of its symbols, the finite reduct
has a finite characteristic sentence: existentially enumerate its
\(k\) different elements, state that every element is one of them,
and give the complete relation table and constant assignments
for \(L_0\).
This sentence is true in \(\mathcal A\), hence in every other
\(S\)-model. There is therefore a bijection preserving each
specified finite \(L_0\).
The set of all bijections between the two \(k\)-element domains
is finite. The subsets preserving each individual symbol have
the finite intersection property just proved, so their full
intersection is nonempty. A member is an isomorphism for the
whole, possibly infinite, signature.
Every model is thus isomorphic to \(\mathcal A\), making it prime.

For any finite tuple, its complete type is one of only finitely
many types realized in \(\mathcal A\). For each different such
type choose one formula distinguishing it from the chosen tuple;
their finite conjunction isolates that tuple's type.
All models are isomorphic, so this isolator works over \(S\).
This proves atomicity, including the unique zero-type.
For equal-type tuples, apply the finite-bijection argument to
the structures with those entries named. It gives an automorphism
with that prescribed finite correspondence, proving homogeneity.
For any parameter set within \(\mathcal A\), which is itself finite,
a consistent complete finite-tuple type has a realization in a
model of its full named theory by compactness.
That model is finite of size \(k\); the same named isomorphism
argument transfers the realizing tuple back to \(\mathcal A\).
This proves the saturation assertion at every parameter scope
available within this finite domain.

In infinite pure equality, finite tuples have equality partitions.
Two tuples with the same partition extend to match a new entry
in either direction, choosing outside the old finite set when
necessary. Formula induction transfers all formulas, proving
completeness and isolation by the equality pattern in every
infinite model, regardless of its size.
Thus all are atomic. A countable such model is prime by the
proved countable package, whereas an uncountable one cannot
inject, hence cannot elementarily embed, into the countable model.

For the incomplete pure-equality one-or-two-element theory, the one-element
model's tuples are isolated by the one-element cardinality
sentence and their forced equalities, so its formula-algebra
atomicity is immediate. But \(\forall y(y=v_0)\) is true at its
element and false at each element of a two-element model.
An elementary embedding is impossible.
This proves the displayed completeness boundary.
\end{proof}
\end{annexitem}

\begin{annexitem}{annex:p3.vaught.functional-bridge}{thm:p03-vaught-functional-bridge}{The graph translation and both inverse directions}
\begin{proof}
Choose a fresh relation \(G_f\) for each positive-arity function,
and impose
\(\forall\bar x\exists!y\,G_f(\bar x,y)\).
In a function model, interpret it by \(f(\bar x)=y\).
Conversely these total single-valued relations define a unique
function at every argument tuple on the same domain.
Relations and individual constants are unchanged.
These two model operations are mutually inverse and preserve
domain cardinality.

For each term \(t(\bar x)\), define a graph evaluation formula
\(E_t(\bar x,z)\). A variable gives \(z=x_i\), a constant gives
\(z=d\), and \(f(t_1,\ldots,t_r)\) gives
\[
 \exists z_1\cdots z_r
   \left(\bigwedge_i E_{t_i}(\bar x,z_i)
                    \land G_f(z_1,\ldots,z_r,z)\right).
\]
Use fresh bound variables at every occurrence.
Induction on the term proves existence and uniqueness of the
selected \(z\) in every graph model, and identifies it with
the function term's value in the corresponding function model.
Translate an atomic relation on terms by existentially selecting
their graph values and applying the unchanged relation.
Translate term equality by selecting the two values and equating
them. Translate Boolean operations and quantifiers compositionally.
The term induction gives atomic truth preservation; formula
induction gives the full truth preservation and retains exactly
the free variables.

In the other direction replace a graph atomic formula
\(G_f(\bar x,y)\) by \(f(\bar x)=y\).
The unique term-value lemma shows that both successive translations
are equivalent to the starting formula in their respective theories.
For a function theory \(S\), let \(S^{\rm gr}\) contain the
translated axioms and all graph axioms.
Models, satisfaction of additional formula constraints and
consistency therefore correspond in both directions.
Complete types correspond by translating every formula and using
these inverse equivalences; partial types and their realization
constraints correspond as well.
If a graph-language consistent formula generated a translated
partial type, its inverse function formula would be a consistent
generator for the original partial type. The forward argument
is the same. This proves principality preservation, and hence
isolation and atomicity preservation for complete types.

A map is elementary in one language exactly when it is in the
other, by full truth preservation. In particular preservation
of graph relations includes preservation of function values.
Elementary substructures of a total graph model remain closed
under each function, since the existential graph witness for
an old argument tuple must occur in the elementary substructure.
The model bijection and map correspondence identify primeness.
The type correspondence over the same named parameter sets
identifies saturation at every specified parameter bound.
Automorphisms and equal-type finite tuples correspond, giving
homogeneity. Finally a countable function list contributes only
countably many finite-arity graph symbols and axioms.
The relational source results therefore transport with every
stated language and cardinal condition retained.
\end{proof}
\end{annexitem}

\begin{annexitem}{annex:p3.vaught.partial-omission}{thm:p03-vaught-partial-omission}{Every partial-type refinement, including the empty tuple}
\begin{proof}
First use a relational/constant countable language.
Run the same fresh witness enumeration and constant-tuple
task enumeration for \(S\), with one task for each listed
partial type and each syntactically distinct tuple of its arity.
For arity zero there is just the empty tuple task.
If the nonempty task list is finite, repeat its entries to obtain
an infinite schedule paired with every witness stage.
Maintain consistency of the finite accumulated \(\Gamma_s\).
At a task for \(p_j\), form the consistent witness extension
\(\gamma_s\), replace its new constant symbols by distinct
free variables, and existentially project its other slots
to obtain the consistent formula \(\beta_s\) in the task's
free-variable slots. This projection argument did not use
completeness of the starting theory.

There must be a \(\varphi\in p_j\) such that
\(\beta_s\land\neg\varphi\) is consistent with \(S\).
If none existed, the consistent \(\beta_s\) would imply every
member of \(p_j\) over \(S\) and would be a principal generator
under the exact definition. Nonprincipality rules this out.
Choose such a \(\varphi\). The existential projection provides
interpretations of the private constant slots witnessing
\(\gamma_s\land\neg\varphi(\bar d)\), so adding this negated
instance preserves consistency.
Retain the witness axiom and all preceding sentences.
For a zero-length tuple, \(\beta_s\) and \(\varphi\) are sentences;
the same consistency and generator argument applies without
any tuple variables or names.

Every finite subset of all stage sentences is satisfiable.
Compactness and the complete named-model lemma give an at most
countable model of them, elementary in the witness-theory model.
Every positive-length tuple has syntactically distinct names
by the retained infinite-alias argument, so at its task for
\(p_j\) a member of that type was made false.
The empty tuple was covered directly at its task.
Thus all listed types are omitted. The empty family just runs
the witness construction. If \(S\) forces infinitude the named
model is infinite; otherwise finite domains are permitted.

For a countable function language, apply the graph bridge to
\(S\) and to the listed partial types.
Their translated types remain consistent and nonprincipal
by the proved inverse generator correspondence.
The relational construction supplies a graph model omitting
them, whose corresponding function model, on the same countable domain,
omits the original types by truth preservation.
This proves the broader interface in full, without identifying
it as Vaught's separately stated atom/atomless theorem.
\end{proof}
\end{annexitem}

\begin{annexitem}{annex:p3.vaught.stability-boundary}{thm:p03-vaught-stability-boundary}{Counting every finite arity from the countable-parameter premise}
\begin{proof}
Induct on finite arity. The zero-type space has one point.
Suppose \(\mathcal P_n(T)\) is countable.
For every \(p\) in it introduce a fresh named tuple and impose
the complete \(p\) at that tuple. These countably many requests
are jointly consistent. Indeed a finite portion requests finitely
many consistent finite conjunctions; completeness of \(T\) makes
each existential closure true in every \(T\)-model, so one such
model supplies a witness tuple for every conjunction.
Compactness gives a common model \(\mathcal M\) and one tuple
\(\bar a_p\) of type \(p\) for each request.
Let \(A\) be the union of their entries; it is countable.

For each \(Q\in\mathcal P_{n+1}(T)\) restricting to \(p\), its
last-coordinate type over \(\bar a_p\) is consistent with the
complete theory naming all of \(A\).
For any finite \(Q\)-conjunction, its last-variable existential
closure belongs to \(p\), so the actual \(\bar a_p\) has a witness
in \(\mathcal M\). That expansion already satisfies every true
sentence over \(A\), proving finite consistency together with
the complete named theory. Compactness supplies a complete
one-type \(q\) over \(A\) extending these constraints.
\(\omega\)-stability makes the space of all such \(q\) countable.

Choose one such extension for each \(Q\), using the stated ambient
Choice convention. The pair \((p,q)\) determines \(Q\): restriction
of \(q\) to the names \(\bar a_p\) determines exactly every
concatenated formula, by the already proved finite-parameter
type correspondence. Two different \(Q\) therefore cannot give
the same pair. This injects \(\mathcal P_{n+1}(T)\) into a product
of two countable sets. It is countable, completing the induction.
Every finite empty-parameter type space is thus countable,
which is precisely smallness.
\end{proof}
\end{annexitem}

\begin{annexitem}{annex:p3.vaught.finite-formulas-categoricity}{thm:p03-vaught-finite-formulas-categoricity}{The actual finite-algebra and two-model argument}
\begin{proof}
For a finite Boolean algebra, the previously proved finite-atom
argument makes every complete filter principal and leaves only
finitely many such filters, one per atom.
Conversely, if all complete filters of a countable formula algebra
are principal, the nonprincipal-algebra lemma rules out an infinite
algebra. Thus clauses 2 and 4 are equivalent, and 2 implies 3.
If there are only finitely many filters, distinct algebra elements
give distinct subsets of that finite set: a nonzero symmetric
difference is realized by a tuple in a \(T\)-model, whose truth
type distinguishes the two elements.
There are at most \(2^k\) subsets of a \(k\)-element set.
The algebra is therefore finite, proving \(3\Rightarrow2\).

Under clause 2, every tuple in any \(T\)-model realizes a principal
complete type and thus an atom. Every model is atomic.
A countably infinite model exists by the naming construction,
and any two such models are isomorphic by full atomic
back-and-forth. This gives \(2\Rightarrow1\).

If clause 2 fails at some arity \(n\), the nonprincipal-algebra
lemma gives a nonprincipal complete \(P\).
The consistent named theory \(\overline P\) has an at most countable
model by the naming construction. Its reduct is infinite, hence
countably infinite, and realizes \(P\).
The simultaneous omission theorem gives a countably infinite
model of \(T\) omitting \(P\).
These models are not isomorphic: isomorphisms preserve all tuple
formulas and hence realization of this complete type.
Thus clause 1 fails. This proves \(1\Rightarrow2\) by contraposition,
and completes every equivalence from the constructions proved here.
\end{proof}
\end{annexitem}
